\chapter{Introducción}

El actual aumento de la infraestructura digital ha ampliado significativamente la superficie de ataque disponible para actores maliciosos, convirtiendo la ciberseguridad en una preocupación crítica para las organizaciones. Durante el periodo 2023–2024, se registró un aumento sin precedentes en la sofisticación de las amenazas; según el Microsoft Digital Defense Report \cite{MS-data}, este incremento ha convertido a las "instituciones gubernamentales y los servicios financieros en los objetivos primordiales, representando la mayor parte de las operaciones de espionaje y ataques de ransomware a nivel global". El vector predominante combina ingeniería social con malware diseñado para evadir controles técnicos y permanecer sin detección por períodos prolongados dentro de los sistemas comprometidos. En este contexto, los registros de eventos de seguridad, comúnmente denominados logs, constituyen la principal fuente de evidencia disponible para detectar, investigar y responder a incidentes. Una infraestructura tecnológica moderna puede generar cientos de miles de eventos por hora, provenientes de fuentes heterogéneas que incluyen sistemas operativos, firewalls, aplicaciones web, contenedores y servicios en la nube.


Sin embargo, el análisis efectivo de logs de seguridad enfrenta desafíos estructurales que las soluciones actuales no han logrado resolver de manera integral. Las plataformas SIEM (\textit{Security Information and Event Management}) representan el estándar vigente para la centralización y correlación de logs; no obstante, dependen fundamentalmente de reglas predefinidas y firmas estáticas para identificar comportamientos maliciosos. Esta dependencia genera dos problemas críticos: por una parte, una alta tasa de falsos positivos que sobrecarga a los analistas; por otra, la incapacidad de detectar ataques que no coincidan exactamente con los patrones preconfigurados.

Más aún, el desafío fundamental no radica en detectar eventos anómalos de manera aislada, sino en comprender cómo una secuencia de eventos dispersos configura la narrativa coherente de un ataque: cuál fue el vector de entrada, qué acciones ejecutó el atacante, qué recursos comprometió y cuál fue el impacto real. Esta tarea de reconstrucción de incidentes requiere razonamiento causal sobre series temporales de eventos heterogéneos, una capacidad de la que carecen los sistemas basados en reglas.

En consecuencia, la carga investigativa recae en analistas humanos expertos; sin embargo, esto representa uno de los principales obstáculos para la madurez de las organizaciones en materia de seguridad, afectando especialmente a las PyMEs que no pueden costear equipos especializados.

La comunidad investigadora ha explorado diversas aproximaciones para superar estas limitaciones. Una primera generación de soluciones incorporó técnicas clásicas de aprendizaje automático para la detección de anomalías en logs, logrando mejoras respecto a los enfoques basados en firmas, aunque manteniendo sensibilidad frente al fenómeno de \textit{log drift}, el cual degrada progresivamente su desempeño en entornos productivos dinámicos. Posteriormente, una segunda generación de propuestas ha comenzado a integrar Modelos de Lenguaje de Gran Escala (LLMs) al análisis de logs, aprovechando su capacidad para comprender el significado semántico del texto más allá de su estructura léxica superficial. En este sentido, estudios recientes destacan que los LLMs permiten capturar relaciones contextuales complejas y patrones semánticos latentes presentes en los mensajes de log, superando limitaciones propias de métodos tradicionales de representación textual como TF-IDF o Word2Vec \cite{cui2025logeval,Akhtar2025}.

Trabajos recientes como LogEval \cite{cui2025logeval}, el estudio de Akhtar et al. \cite{Akhtar2025} y propuestas aplicadas como LogAn \cite{Gupta2025} han evidenciado que los modelos de lenguaje de gran escala poseen un alto potencial para superar las limitaciones de los enfoques tradicionales de aprendizaje automático en tareas de análisis de logs, particularmente en procesos de parseo, detección de anomalías, diagnóstico y generación de resúmenes. En particular, LogEval reporta mejoras relevantes en distintas tareas de análisis operacional al evaluar múltiples modelos sobre benchmarks especializados \cite{cui2025logeval}. En el ámbito de los centros de operaciones de seguridad (\textit{SOC}), investigaciones recientes también han sistematizado el potencial de los LLMs para reducir la fatiga de alertas, automatizar procesos de triage y mejorar la interpretación de eventos de seguridad \cite{Akhtar2025}. Sin embargo, estos trabajos comparten una limitación arquitectónica fundamental: operan principalmente como pipelines lineales y pasivos que procesan eventos de manera relativamente aislada y entregan resultados estáticos, sin incorporar ciclos de razonamiento iterativo que permitan investigar activamente un incidente. En esta línea, LogAn constituye la propuesta más cercana en términos arquitectónicos, al integrar inferencia causal aplicada al dominio de IT Operations; no obstante, su alcance no contempla el razonamiento sobre tácticas adversariales ni la reconstrucción integral de incidentes de ciberseguridad \cite{Gupta2025}.

El presente trabajo documenta las etapas de diseño y evaluación de un agente basado en LLMs para el análisis semántico de incidentes de seguridad a partir de logs. A diferencia de los pipelines convencionales, el sistema implementa un ciclo de razonamiento agentivo sustentado en el patrón \textit{Plan-and-Execute}, mediante el cual un módulo Planificador construye un plan ordenado de análisis sobre un resumen estructurado de los logs, y un módulo Ejecutor lo implementa paso a paso invocando agentes LLM especializados y herramientas determinísticas: consulta de reputación de IPs, recuperación de vulnerabilidades CVE, generación de comandos de mitigación validados y síntesis del diagnóstico final. La contribución técnica combina detección semántica de anomalías mediante embeddings derivados de LLMs, razonamiento causal iterativo y recuperación aumentada de contexto mediante técnicas RAG (\textit{Retrieval-Augmented Generation}) sobre fuentes de inteligencia de amenazas. Desde una perspectiva aplicada, la solución se orienta a organizaciones que no disponen de equipos especializados, particularmente pequeñas y medianas empresas del contexto chileno y latinoamericano. Entre los resultados esperados se incluyen la definición formal de la arquitectura del agente, el desarrollo de un prototipo funcional evaluado sobre datasets públicos de referencia y el establecimiento de métricas para medir la calidad de la reconstrucción de incidentes, dimensión escasamente abordada en la literatura actual y que abre una línea de investigación propia.